\documentclass[11pt,a4paper]{article}
\usepackage[margin=1in]{geometry}
\usepackage{amsmath,amssymb,amsthm}
\usepackage{mathtools}
\usepackage{enumitem}

\newtheorem{theorem}{Theorem}[section]
\newtheorem{proposition}[theorem]{Proposition}
\newtheorem{lemma}[theorem]{Lemma}
\newtheorem{corollary}[theorem]{Corollary}
\theoremstyle{definition}
\newtheorem{definition}[theorem]{Definition}
\newtheorem{axiom}{Axiom}
\theoremstyle{remark}
\newtheorem{remark}[theorem]{Remark}
\newtheorem*{remark*}{Remark}

\newcommand{\CP}{\mathbb{CP}}
\newcommand{\C}{\mathbb{C}}
\newcommand{\R}{\mathbb{R}}
\newcommand{\Z}{\mathbb{Z}}
\newcommand{\Sym}{\mathrm{Sym}}
\newcommand{\End}{\mathrm{End}}
\newcommand{\SL}{\mathrm{SL}}
\newcommand{\GL}{\mathrm{GL}}
\newcommand{\PGL}{\mathrm{PGL}}
\newcommand{\SO}{\mathrm{SO}}

\title{The Minimum Self-Consistent Relational Structure}
\author{J.\,R.~Manuel}
\date{April 2026}

\begin{document}
\maketitle

%% ====================================================================
\section{Definitions}\label{sec:definitions}
%% ====================================================================

\begin{definition}[Primitive space]\label{def:primitive}
A \emph{primitive space} is a finite-dimensional complex vector space $S$ with
$\dim_\C(S)\geq 2$.
The elements of $S$ are called \emph{primitive entities}.
The condition $\dim(S)\geq 2$ encodes the requirement that at least two
distinguishable entities exist: below two, there is no relation, hence no
structure.
\end{definition}

\begin{definition}[Relational structure]\label{def:relational}
A \emph{relational structure} on a primitive space $S$ is a triple
$(S,V,\pi)$ where:
\begin{enumerate}[nosep]
\item $V$ is a finite-dimensional complex vector space (the \emph{state space}),
\item $\pi\colon S\times S \to V$ is the canonical bilinear map
$\pi(u,v)=u\otimes v$, identifying $V$ with $S\otimes S$.
\end{enumerate}
The bilinear map $\pi$ encodes the assertion that every state arises from
the interaction (tensor product) of primitive entities.
\end{definition}

\begin{definition}[Dynamical and substrate decomposition]\label{def:decomp}
The tensor product $S\otimes S$ decomposes under the symmetric group
$\mathfrak{S}_2$ as
\[
S\otimes S \;=\; \Sym^2(S)\;\oplus\;\Lambda^2(S).
\]
The \emph{dynamical algebra} is $\Sym^2(S)$: the symmetric interactions,
which carry the $\SL(S)$-adjoint representation and encode causal
structure (see Corollary~\ref{cor:sl2}).
The \emph{substrate} is $\Lambda^2(S)$: the antisymmetric part, which
carries the invariant volume forms and provides the medium in which
relations are embedded.
\end{definition}

\begin{definition}[Self-consistency]\label{def:selfconsist}
A relational structure $(S,V,\pi)$ is \emph{self-consistent} if
$V=S\otimes S$---that is, the state space is neither more nor less than
the bilinear interactions of the primitives.
\end{definition}

\begin{remark}[Why $V=S\otimes S$]\label{rem:why_tensor}
Self-consistency demands that the state space be generated by the
primitive entities under their mutual interaction.
The tensor product $S\otimes S$ is the universal bilinear construction:
it is the smallest space through which every bilinear map on $S$ factors.
To take $V\subsetneq S\otimes S$ would discard bilinear interactions,
breaking universality.
To take $V\supsetneq S\otimes S$ would require structure not generated
by the primitives, violating self-containment.
\end{remark}

%% ====================================================================
\section{Axioms}\label{sec:axioms}
%% ====================================================================

\begin{axiom}[Distinguishability]\label{ax:distinguish}
The primitive space $S$ is a finite-dimensional complex vector space with
$\dim_\C(S)\geq 2$.
\end{axiom}

\begin{axiom}[Self-consistency]\label{ax:selfconsist}
The state space is $V=S\otimes S$.
\end{axiom}

\begin{axiom}[Unique substrate]\label{ax:substrate}
$\dim_\C(\Lambda^2(S))=1$.
\end{axiom}

\begin{remark}[On the axiom count]\label{rem:axiom_count}
These three axioms are not independent.
Axiom~\ref{ax:substrate} alone forces $\dim(S)=2$, which then satisfies
Axiom~\ref{ax:distinguish}.
The logical content is:
\begin{enumerate}[nosep]
\item $V=S\otimes S$ (self-consistency),
\item $\dim(\Lambda^2(S))=1$ (unique substrate),
\end{enumerate}
from which everything follows.
Axiom~\ref{ax:distinguish} is retained for conceptual clarity: it makes
explicit the minimal relational requirement before the substrate
constraint is imposed.
\end{remark}

%% ====================================================================
\section{The Main Theorem}\label{sec:main}
%% ====================================================================

\begin{theorem}[Minimum self-consistent relational structure]\label{thm:main}
Under Axioms~\ref{ax:distinguish}--\ref{ax:substrate}:
\begin{enumerate}[label=(\roman*)]
\item $\dim(S)=2$, so $S\cong\C^2$.
\item $\dim(V)=4$, so $V\cong\C^4$.
\item Setting $H\coloneqq\dim(\Sym^2(S))$, we have $H=3$ and the
self-consistency identity
\begin{equation}\label{eq:selfconsist}
(H-1)^2 = H+1
\end{equation}
is satisfied.
\end{enumerate}
\end{theorem}

\begin{proof}
Let $n=\dim_\C(S)$.

\medskip\noindent\textbf{(i)} The exterior power $\Lambda^2(S)$ has dimension
$\binom{n}{2}=n(n-1)/2$.
Axiom~\ref{ax:substrate} requires $n(n-1)/2=1$, i.e., $n(n-1)=2$.
This is a quadratic in $n$ with roots $n=2$ and $n=-1$.
Since $n$ is a positive integer, $n=2$.

\medskip\noindent\textbf{(ii)} By Axiom~\ref{ax:selfconsist},
$\dim(V)=\dim(S\otimes S)=n^2=4$.

\medskip\noindent\textbf{(iii)} $\dim(\Sym^2(S))=\binom{n+1}{2}=n(n+1)/2$.
At $n=2$: $H=2\cdot 3/2=3$.
The decomposition $\dim(V)=\dim(\Sym^2(S))+\dim(\Lambda^2(S))$ gives
$4=3+1$, which is $H+1=4$.
Meanwhile, $(H-1)^2=(3-1)^2=4$.
Hence $(H-1)^2=H+1$.

\medskip\noindent
Alternatively, equation~\eqref{eq:selfconsist} is derived from the
dimension identity. We have $\dim(V)=n^2$ and
$\dim(\Sym^2(S))=n(n+1)/2=H$, $\dim(\Lambda^2(S))=n(n-1)/2$.
The decomposition $n^2=H+n(n-1)/2$ gives
$H=n^2-n(n-1)/2=n(n+1)/2$.
The constraint $\dim(\Lambda^2(S))=1$ means $n(n-1)/2=1$, so
$\dim(V)=H+1$.
Now $(H-1)^2=\bigl(\tfrac{n(n+1)}{2}-1\bigr)^2$.
At $n=2$ (the unique solution): $(3-1)^2=4=3+1=H+1$.
To see that this is the \emph{content} of the unique substrate axiom in
the self-consistency equation: substituting $\dim(V)=H+1$ into the
identity $\dim(V)=(H-1)^2$ would require $(H-1)^2=H+1$, i.e.,
$H^2-3H=0$, giving $H=3$ as the unique positive solution.
\end{proof}

\begin{remark}[The equation $(H-1)^2=H+1$]\label{rem:equation}
In the present context, equation~\eqref{eq:selfconsist} is a
consequence of the dimension identity $\dim(V)=H+1$ forced by
$\dim(\Lambda^2(S))=1$.
It acquires a second, independent meaning when read in the tangent
bundle setting (Corollary~\ref{cor:quartic}): the left side $(H-1)^2$
counts the dimension of traceless symmetric bilinear forms on $H-1$
independent directions---the number of independent bilinear products
the dynamical algebra can produce---and the right side $H+1=\dim(V)$ is
the dimension of the full state space.
The equation then says: the self-interaction of the dynamical algebra
\emph{exactly fills} the state space.
That the same equation arises from two different starting points
(unique substrate, and tangent bundle self-consistency) is the
mathematical content of the uniqueness.
\end{remark}

%% ====================================================================
\section{Corollaries}\label{sec:corollaries}
%% ====================================================================

\begin{corollary}[Dynamical algebra is $\mathfrak{sl}(2,\C)$]\label{cor:sl2}
$\Sym^2(\C^2)$ is isomorphic, as an $\SL(2,\C)$-representation, to the
adjoint representation $\mathfrak{sl}(2,\C)$.
\end{corollary}

\begin{proof}
The natural $\SL(2,\C)$-action on $S=\C^2$ induces an action on
$\Sym^2(S)$.
The adjoint representation of $\SL(2,\C)$ on its Lie algebra
$\mathfrak{sl}(2,\C)$ has dimension~$3$.
Both are irreducible representations of dimension~$3$ of $\SL(2,\C)$
(the spin-$1$ representation), hence isomorphic.

Explicitly: let $\{e_0,e_1\}$ be a basis of $S$.
The symmetric products $e_0^2,\,e_0 e_1,\,e_1^2$ form a basis of
$\Sym^2(S)$.
Under the standard identification with homogeneous polynomials of
degree~$2$ in two variables, the $\SL(2,\C)$-action by linear
substitution gives the spin-$1$ (adjoint) representation.
\end{proof}

\begin{corollary}[Substrate is the symplectic form]\label{cor:symplectic}
$\Lambda^2(\C^2)\cong\C$ carries the unique (up to scale)
$\SL(2,\C)$-invariant antisymmetric bilinear form
$\varepsilon\colon S\times S\to\C$.
This is the spinor metric (the $\varepsilon$-spinor of Penrose--Rindler).
\end{corollary}

\begin{proof}
$\Lambda^2(\C^2)$ is one-dimensional, hence carries the trivial
$\SL(2,\C)$-representation (since $\SL(2,\C)$ preserves the volume form
by definition).
The generator $e_0\wedge e_1$ defines $\varepsilon_{AB}$ with
$\varepsilon_{01}=-\varepsilon_{10}=1$, the standard symplectic form on
$\C^2$.
That $\SL(2,\C)$ is the group preserving $\varepsilon$ (up to sign)
confirms the identification.
\end{proof}

\begin{corollary}[Twistor space]\label{cor:twistor}
The projectivization of the state space is $\mathbb{P}(V)=\CP^3$, which
is Penrose's twistor space for conformally compactified Minkowski space.
\end{corollary}

\begin{proof}
$V=\C^4$, so $\mathbb{P}(V)=\CP^3$.
The identification with Penrose's twistor space $\mathbb{PT}$ follows
from the $\SL(2,\C)$-decomposition.
The state space $V=S\otimes S=\C^2\otimes\C^2$ carries the natural
$\SL(2,\C)\times\SL(2,\C)$ action.
In the twistor formalism, a twistor $Z^\alpha=(\omega^A,\pi_{A'})$
decomposes under the spin group $\SL(2,\C)_L\times\SL(2,\C)_R$ of the
conformal group, with $\omega^A\in S_L$ and $\pi_{A'}\in S_R$.
The conformal group of compactified Minkowski space is
$\SL(4,\C)/\Z_4\cong\SO(6,\C)$, acting naturally on $\C^4$.
The twistor fibration $\CP^3\to S^4$ is the standard Penrose fibration.
\end{proof}

\begin{corollary}[Holomorphic sections of the tangent bundle]\label{cor:sections}
$\Sym^2(\C^2)\cong H^0(\CP^1,\mathcal{O}(2))$: the dynamical algebra
is the space of holomorphic sections of the tangent bundle of the
Riemann sphere.
\end{corollary}

\begin{proof}
The tangent bundle of $\CP^1$ is the line bundle $\mathcal{O}(2)$.
Its holomorphic sections are homogeneous polynomials of degree~$2$ in
the homogeneous coordinates $[\pi_0:\pi_1]$ of $\CP^1$:
\[
\eta(\pi) = s_0\,\pi_0^2 + s_1\,\pi_0\pi_1 + s_2\,\pi_1^2,
\qquad (s_0,s_1,s_2)\in\C^3.
\]
This is precisely $\Sym^2(S^\vee)\cong\Sym^2(S)$ (using the symplectic
form $\varepsilon$ to identify $S\cong S^\vee$).
Hence $\dim H^0(\CP^1,\mathcal{O}(2))=3=H$.
\end{proof}

\begin{corollary}[The quartic factorization]\label{cor:quartic}
$\CP^1$ is the unique compact complex projective space $\CP^n$
($n\geq 1$) whose tangent bundle satisfies the self-consistency
equation.
\end{corollary}

\begin{proof}
For $\CP^n$, $\dim H^0(T\CP^n)=(n+1)^2-1$ (the generators of
$\PGL(n+1,\C)$).
Substituting $H=(n+1)^2-1$ into $(H-1)^2=H+1$:
\[
((n+1)^2-2)^2 = (n+1)^2.
\]
Expanding: $n^4+4n^3+n^2-6n=0$, which factors completely over $\Z$:
\[
n(n-1)(n+2)(n+3)=0.
\]
Roots: $n\in\{0,1,-2,-3\}$. The unique positive integer is $n=1$.
\end{proof}

%% ====================================================================
\section{Critical Assessment}\label{sec:critical}
%% ====================================================================

\subsection{Status of Axiom~\ref{ax:selfconsist}: $V=S\otimes S$}

This is the strongest axiom and the one most open to challenge.
The claim is that the state space equals the full tensor product of the
primitive space with itself.
Three points bear scrutiny.

\textbf{Why not $V=\Sym^2(S)$?}
One might argue that only the symmetric part---the ``observable''
interactions---should constitute the state space, discarding the
antisymmetric part as unphysical.
This would give $V=\Sym^2(S)$, $\dim(V)=n(n+1)/2$, and no substrate
at all.
The counter-argument is structural: the tensor product $S\otimes S$ is
the \emph{universal} bilinear construction.
Every bilinear map $S\times S\to W$ factors through $S\otimes S$, not
through $\Sym^2(S)$.
To restrict to $\Sym^2(S)$ is to impose a symmetry condition
\emph{before} asking what the structure requires---it is an additional
axiom, not a simplification.
The full tensor product is the canonical, assumption-free choice.
Furthermore, discarding $\Lambda^2(S)$ eliminates the symplectic
structure that defines the spinor metric, destroying the very object
(the $\varepsilon$-spinor) that makes $S$ a spin space.

\textbf{Why not $V=S\otimes S'$ for $S'\neq S$?}
Self-consistency requires that the state space be generated by
interactions of the \emph{same} primitive entities with themselves.
Introducing a second, independent primitive space $S'$ would add
external structure not determined by $S$, violating the closure
condition.

\textbf{Why not higher tensor powers?}
One could consider $V=S^{\otimes k}$ for $k\geq 3$.
The bilinear requirement ($k=2$) is the minimum: a single interaction
between two entities.
Higher tensor powers describe multi-party interactions and are
\emph{derived} from the bilinear structure, not primitive.
The tensor product $S\otimes S$ is the irreducible case.

\subsection{Status of Axiom~\ref{ax:substrate}: $\dim(\Lambda^2(S))=1$}

This axiom encodes uniqueness of the substrate.
It is the sharpest constraint and admits two readings:

\textbf{Reading 1: Minimality.}
The substrate is the medium through which relations are embedded.
Multiple independent substrates ($\dim(\Lambda^2(S))>1$) would require
a meta-substrate to relate them, producing infinite regress.
One substrate is the minimum that avoids this regress.
Under this reading, Axiom~\ref{ax:substrate} is a minimality condition,
not a logical necessity.

\textbf{Reading 2: Non-degeneracy.}
For $S$ to be a spin space, it must carry a non-degenerate symplectic
form $\varepsilon\in\Lambda^2(S^\vee)$.
Non-degeneracy requires $\Lambda^2(S)$ to be one-dimensional (so that
the symplectic form is unique up to scale and cannot degenerate into
a lower-rank object).
For $\dim(S)\geq 3$, $\Lambda^2(S)$ has dimension $\geq 3$, and the
``symplectic form'' is no longer unique---there is a moduli space of
symplectic structures.
The condition $\dim(\Lambda^2(S))=1$ is equivalent to requiring that the
antisymmetric structure be \emph{rigid}: no continuous deformations, no
moduli, no choices.

Both readings lead to the same constraint, but neither is a logical
theorem---each rests on an interpretive commitment (no regress, or
rigidity of the antisymmetric structure).
This is the weakest point of the axiom system.

\subsection{Alternative axiom systems}

\textbf{Dropping Axiom~\ref{ax:substrate}.}
Without the unique substrate condition, $\dim(S)$ is unconstrained
(subject to $\dim(S)\geq 2$).
For $\dim(S)=3$: $V=\C^9$, $\Sym^2(S)=\C^6$, $\Lambda^2(S)=\C^3$.
The self-consistency equation $(H-1)^2=H+1$ with $H=6$ gives
$25\neq 7$: it fails.
For general $n$, setting $H=n(n+1)/2$ and $\dim(V)=n^2$, the condition
$(H-1)^2=\dim(V)$ becomes
$\bigl(\tfrac{n(n+1)}{2}-1\bigr)^2=n^2$, which gives $n=2$ as the
unique positive integer solution with $n\geq 2$.
Taking the positive square root: $\tfrac{n(n+1)}{2}-1=n$, i.e.,
$n^2+n-2=2n$, i.e., $n^2-n-2=0=(n-2)(n+1)=0$, giving $n=2$.
(The negative root $\tfrac{n(n+1)}{2}-1=-n$ gives $n^2+3n-2=0$, which
has no positive integer solution.)
Direct verification: at $n=2$,
$\bigl(\tfrac{2\cdot 3}{2}-1\bigr)^2=(3-1)^2=4=2^2$, confirmed.
This shows that even replacing Axiom~\ref{ax:substrate} with the
self-consistency equation $(H-1)^2=\dim(V)$ directly, $n=2$ is the
unique solution with $n\geq 2$.
In other words, the self-consistency equation and the unique substrate
condition are \emph{equivalent} in this setting.

\textbf{Real instead of complex.}
If $S$ is a real vector space, $\Lambda^2(S)$ carries an orthogonal
(not symplectic) structure.
The condition $\dim(\Lambda^2(S))=1$ still gives $\dim(S)=2$, yielding
$V=\R^4$ and $H=3$.
But the dynamical algebra is $\mathfrak{sl}(2,\R)$, the split real
form, whose Killing form has signature $(1,2)$.
The absence of complex structure means $\R^4$ does not projectivize to
a complex manifold, and the holomorphic geometry essential to twistor
theory is unavailable.
The complex structure on $S$ is necessary for the correct signature and
for the holomorphic geometry of twistor theory.
This is consistent with Theorem~9 of the main paper (uniqueness of
$\C$ among normed division algebras with an ordering-sensitive
encoding).

\textbf{Quaternionic.}
If $S$ is a right $\mathbb{H}$-module with $\dim_\mathbb{H}(S)=1$
(the quaternionic analogue of $\C^2$), then
$S\otimes_\mathbb{H}S$ is not well-defined in the same sense (the
tensor product over $\mathbb{H}$ does not yield a free
$\mathbb{H}$-module in general because $\mathbb{H}$ is
non-commutative).
Treating $S$ as $\C^4$ with a quaternionic structure $j$, the
self-consistency analysis leads to $\dim(\Lambda^2_\C(S))=6$, violating
the unique substrate condition.

\subsection{Summary of logical status}

\begin{center}
\renewcommand{\arraystretch}{1.3}
\begin{tabular}{lll}
\textbf{Axiom} & \textbf{Status} & \textbf{Weakest point} \\
\hline
A1: $\dim(S)\geq 2$ & Definitional &
``Two entities'' is a starting premise \\
A2: $V=S\otimes S$ & Structural &
Why not a subspace? (Answer: universality) \\
A3: $\dim(\Lambda^2(S))=1$ & Minimality &
``One substrate'' requires interpretation \\
\end{tabular}
\end{center}

\medskip\noindent
The strongest result is that A2 and A3 are not independent in their
consequences: the self-consistency equation $(H-1)^2=\dim(V)$ with
$V=S\otimes S$ has $n=2$ as its unique solution with $n\geq 2$, and at
$n=2$ one automatically gets $\dim(\Lambda^2(S))=1$.
Conversely, $\dim(\Lambda^2(S))=1$ forces $n=2$, which then satisfies
the self-consistency equation.
The two constraints are equivalent.
This is either a strength (the axiom system is overdetermined and
consistent) or a weakness (one axiom is redundant), depending on
perspective.

\bigskip
\begin{remark*}[What is proven vs.\ assumed]
\emph{Proven:} Given a complex vector space $S$ with $\dim(S)=2$, the
tensor product $V=S\otimes S=\C^4$ decomposes as
$\Sym^2(S)\oplus\Lambda^2(S)=\C^3\oplus\C$, and this is the unique
dimension satisfying $(H-1)^2=H+1$.
The corollaries (adjoint representation, symplectic form, twistor space,
tangent bundle sections) are standard results in representation theory
and algebraic geometry.

\emph{Assumed:} That the state space of a self-consistent relational
structure must be $S\otimes S$ (not some other functor of $S$), and
that the substrate must be unique.
These are interpretive commitments grounded in universality and
minimality, not logical necessities.
The mathematical content is rigorous; the philosophical motivation for
the axioms is, as with all axiom systems, a matter of judgment.
\end{remark*}

\end{document}
